\exercisegroup

\begin{enumerate}
\def\labelenumi{\arabic{enumi})}
\tightlist
\item
   设 K 是拓扑域 R、C、H 之一。设 V 是具有可数生成系的 K-向量空间，并赋予其作为拓扑向量空间的最细拓扑（EVT，第 I 章，§1，第 11 页，推论 4）。对每个整数 \(n \geqslant 1\)，令 \(\mathrm{B}_{n}(\mathrm{V})\) 为 \(V^{n}\) 中由自由序列组成的子空间，令 \(\mathrm{G}_{n}(\mathrm{V})\) 为 V 的 n 维子空间所成的集合。记 \(p: \mathrm{B}_{n}(\mathrm{V}) \to \mathrm{G}_{n}(\mathrm{V})\) 为把序列 \(v \in \mathrm{B}_{n}(\mathrm{V})\) 对应到其生成的向量子空间的映射，并在集合 \(\mathrm{G}_{n}(\mathrm{V})\) 上赋予使 p 连续的最粗拓扑。
\end{enumerate}

\begin{enumerate}
\def\labelenumi{\alph{enumi})}
\item
   证明 \(\mathbf{GL}(n,\mathrm{K})\) 在 \(\mathrm{B}_{n}(\mathrm{V})\) 上的右作用可由下式定义：\(v \cdot g = (\sum_{i=1}^{n} v_{i} a_{i,j})_{1 \leqslant j \leqslant n}\)，其中 \(v = (v_{1}, \ldots, v_{n}) \in \mathrm{B}_{n}(\mathrm{V})\)，\(g = (a_{i,j}) \in \mathbf{GL}(n, \mathrm{K})\)。
\item
   证明映射 p 在商空间 \(\mathrm{B}_{n}(\mathrm{V})/\mathbf{G}\mathbf{L}(n,\mathrm{K})\) 到 \(\mathrm{G}_{n}(\mathrm{V})\) 上诱导同胚。推出空间 \(\mathrm{G}_{n}(\mathrm{V})\) 是仿紧的；若 V 有限维，则它是可度量化的。
\item
   证明 \(\mathrm{B}_{n}(\mathrm{V})\) 是以 \(\mathbf{GL}(n,\mathrm{K})\) 为群、以 \(\mathbf{G}_{n}(\mathrm{V})\) 为底的主纤维丛。
\item
   假设 V 是无限维的。证明空间 \(\mathrm{B}_{n}(\mathrm{V})\) 是可缩的。推出空间 \(\mathrm{G}_{n}(\mathrm{V})\) 是群 \(\mathbf{GL}(n,\mathrm{K})\) 的分类空间。
\end{enumerate}

\begin{enumerate}
\def\labelenumi{\arabic{enumi})}
\setcounter{enumi}{1}
\tightlist
\item
  设 H 是无限维的分离预希尔伯特空间。令 \((e_{i})_{i\geqslant0}\) 是生成 H 的稠密子空间的正交归一向量族。对每个整数 \(n\geqslant1\)，令 \(V_{n}(H)\) 为 \(H^{n}\) 中由正交归一序列 \((u_{1},\ldots,u_{n})\) 组成的子空间。
\end{enumerate}

\begin{enumerate}
\def\labelenumi{\alph{enumi})}
\item
  证明存在唯一的连续线性映射 T: H → H，使得对每个整数 i 都有 \(\mathrm{T}(e_{i}) = e_{i+1}\)。证明 \(\|T(x)\| = \|x\|\) 对所有 \(x \in H\) 都成立。
\item
  令 H' 为 H 中与 \(e_{i}\) 正交的向量所成的子空间，其中 i \textless{} n。证明以下条件等价：(i) x 与 \(\mathrm{T}^{n}(x)\) 线性相关；(ii) \(\mathrm{T}(x) = x\)；(iii) \(x \in H'\)。
\item
  证明存在唯一的同伦 \(\sigma\colon\mathrm{V}_{n}(\mathrm{H})\times\mathbf{I}\to\mathrm{V}_{n}(\mathrm{H})\)，使得对每个 \((u_{1},\ldots,u_{n})\in\mathrm{V}_{n}(\mathrm{H})\) 和每个 \(t\in\mathbf{I}\)，\(\sigma((u_{1},\ldots,u_{n}),t)\) 是由族 \((1-t)(u_{1},\ldots,u_{n})+t(\mathrm{T}^{n}(u_{1}),\ldots,\mathrm{T}^{n}(u_{n}))\) 经过正交归一化过程得到的序列。
\item
  证明存在唯一的同伦 \(\tau\colon\mathrm{V}_{n}(\mathrm{H}^{\prime})\times\mathbf{I}\to\mathrm{V}_{n}(\mathrm{H})\)，使得对每个 \((u_{1},\ldots,u_{n})\in\mathrm{V}_{n}(\mathrm{H}^{\prime})\) 和每个 \(t\in\mathbf{I}\)，\(\sigma((u_{1},\ldots,u_{n}),t)\) 是由族 \((1-t)(u_{1},\ldots,u_{n})+t(e_{1},\ldots,e_{n})\) 经过正交归一化过程得到的序列。
\item
  证明空间 \(V_{n}(H)\) 是可缩的。
\end{enumerate}

\begin{enumerate}
\def\labelenumi{\arabic{enumi})}
\setcounter{enumi}{2}
\tightlist
\item
  设 H 是分离的预希尔伯特空间。设 n 是整数且 \(\geqslant 1\)；记 \(V_{n}(H)\) 为由正交归一序列组成的 \(H^{n}\) 的子空间。设 G 是 \(O(n, \mathbf{R})\) 的闭子群。
\end{enumerate}

\begin{enumerate}
\def\labelenumi{\alph{enumi})}
\tightlist
\item
   证明正交群 O(n, R)（LIE，第 IX 章，§3，第 5 号，第 22 页）在空间 \(V_{n}(H)\) 上按下式适当且自由地作用：
\end{enumerate}

\[
(u _ {1}, \dots , u _ {n}) \cdot (a _ {i, j}) = \left(\sum_ {i = 1} ^ {n} a _ {i, j} u _ {i}\right).
\]

\begin{enumerate}
\def\labelenumi{\alph{enumi})}
\setcounter{enumi}{1}
\item
  设 G 是 \(O(n,\mathbf{R})\) 的闭子群。证明典范满射使 \(V_{n}(\mathrm{H})\) 成为以 G 为群、以 \(V_{n}(\mathrm{H})/\mathrm{G}\) 为底的主纤维丛（先处理 \(G = O(n,\mathbf{R})\) 的情形）。
\item
  证明商空间 \(V_{n}(H)/G\) 是可度量化的。（在 \(H^{n}\) 上赋予其自然的预希尔伯特空间结构，并考虑 \(V_{n}(H)\) 上由 \(d(u,v)=\inf_{g\in G}\|u-v\cdot g\|\) 给出的规度 d。）
\item
  证明商空间 \(V_{n}(H)/G\) 是 G 的分类空间。
\item
  推出每个紧李群都有一个作为可度量拓扑空间的分类空间。（应用 LIE，第 IX 章，§9，第 2 号，第 91 页，定理 1。）
\end{enumerate}

\begin{enumerate}
\def\labelenumi{\arabic{enumi})}
\setcounter{enumi}{3}
\item
  通过考察复希尔伯特空间，以类似方式构造酉群 U(n, C)（LIE，第 IX 章，§，第 4 号，第 21 页）的可度量分类空间。
\item
  \begin{enumerate}
  \def\labelenumii{\alph{enumii})}
  \tightlist
  \item
    设 X 是仿紧拓扑空间。证明每个以 X 为底、以 R 为群的主纤维丛都是可平凡化的。
  \end{enumerate}
\end{enumerate}

\begin{enumerate}
\def\labelenumi{\alph{enumi})}
\setcounter{enumi}{1}
\tightlist
\item
  设 X 是亚历山德罗夫半直线（TG，第 IV 章，第 49 页，练习 12）。构造一个以 X 为底、以 R 为群且不可平凡化的主纤维丛。
\end{enumerate}

\begin{enumerate}
\def\labelenumi{\arabic{enumi})}
\setcounter{enumi}{5}
\tightlist
\item
  设 X 为 \(R \times \{0,1\}\) 的商空间，商取使点 \((t,0)\) 与 \((t,1)\) 对每个 \(x \in R_{+}^{*}\) 都等价的最细等价关系。记 p: \(R \times \{0,1\} \to X\) 为典范满射。
\end{enumerate}

\begin{enumerate}
\def\labelenumi{\alph{enumi})}
\item
  证明拓扑空间 X 同伦于一点。
\item
  对于哪些点 \(x \in X\)，带基点空间 (X, x) 是可缩的？
\item
  令 \(U = p(\mathbf{R} \times \{0\})\)，\(V = p(\mathbf{R} \times \{\mathrm{V}\})\)。设 G 是拓扑群。构造从连续映射集合 \(\mathcal{C}(\mathbf{R}_{+}^{*}; \mathbf{G})\)（从 \(R_{+}^{*}\) 到 G）到以 X 为底、以 G 为群的主纤维丛同构类集合的双射。
\item
  推出存在以 X 为底、以 R 为群且不可平凡化的主纤维丛。
\end{enumerate}

\begin{enumerate}
\def\labelenumi{\arabic{enumi})}
\setcounter{enumi}{6}
\tightlist
\item
  设 G 是拓扑群，E 是以 \(S_{2}\) 为底、以 G 为群的主纤维丛。记 \(S_{2}^{+}\) 和 \(S_{2}^{-}\) 为 \(S_{2}\) 与分别由条件 \(z \geqslant 0\) 和 \(z \leqslant 0\) 定义的半空间的交，即两个半球；将 \(S_{2}^{+} \cap S_{2}^{-}\) 识别为 \(S_{1}\)。
\end{enumerate}

\begin{enumerate}
\def\labelenumi{\alph{enumi})}
\item
  证明丛 E 有截面 \(s^{+}\)（分别为 \(s^{-}\)），分别定义在 \(S_{2}^{+}\)（分别在 \(S_{2}^{-}\)）上。证明存在唯一连续映射 \(f\colon S_{1}\to G\)，使得 \(s^{+}(x)=s^{-}(x)\cdot f(x)\) 对所有 \(x\in S_{1}\) 都成立。
\item
  证明从 \(S_{1}\) 到 G 的映射 \(x \mapsto c(x)c(1)^{-1}\) 的严格同伦类不依赖于截面 \(s^{+}\) 和 \(s^{-}\) 的选择；将其记为 \(c(\mathrm{E})\)。
\item
  证明，两个主纤维丛 E 和 \(E'\)（底为 \(S_{2}\)，群为 G）同构，当且仅当 \(c(\mathrm{E}) = c(\mathrm{E}')\)。
\item
  证明对每个类 \(c \in \pi_{1}(G, e)\)，存在以 \(S_{2}\) 为底、以 G 为群的主纤维丛 E，使得 \(c(E) = c\)。
\end{enumerate}

\begin{enumerate}
\def\labelenumi{\arabic{enumi})}
\setcounter{enumi}{7}
\tightlist
\item
  设 X 是仿紧拓扑空间，令 S(X) 为其悬垂（I，第 149 页，练习 1）；记 \(p\colon X \times [0,1] \to S(X)\) 为典范满射。令 \(a\) 为 X 中的一点。
\end{enumerate}

\begin{enumerate}
\def\labelenumi{\alph{enumi})}
\item
  设 G 是拓扑群，E 是以 S(X) 为底、以 G 为群的主纤维丛。证明 E 限制到子空间 \(p(\mathrm{X} \times [0,1/2])\) 后有连续截面 \(s_{0}\)，且限制到子空间 \(p(\mathrm{X} \times [1/2,1])\) 后有连续截面 \(s_{1}\)。
\item
  证明存在唯一映射 \(f\colon \mathrm{X}\to \mathrm{G}\)，使得 \(s_1(x) = s_0(x)\cdot f(x)\) 对所有 \(x\in \mathrm{X}\) 都成立。
\item
  证明在 \([(X,a);(G,e)]\) 中，带基点映射 \(x \mapsto f(x)f(a)^{-1}\) 的严格同伦类不依赖于截面 \(s_{0}\) 和 \(s_{1}\) 的选择；将其记为 \(c(\mathrm{E})\)。
\item
  证明，以 S(X) 为底、以 G 为群的两个主纤维丛 E 和 \(E'\) 同构，当且仅当 \(c(\mathrm{E}) = c(\mathrm{E}')\)。
\item
  证明对每个类 \(c \in [(X, a); (G, e)]\)，存在以 S(X) 为底、以 G 为群的主纤维丛 E，使得 \(c(E) = c\)。
\end{enumerate}

\begin{enumerate}
\def\labelenumi{\arabic{enumi})}
\setcounter{enumi}{8}
\tightlist
\item
  我们将复射影直线 \(\mathbf{P}^{1}(\mathbf{C})\) 识别为 \(C^{2}\) 中的直线集。令 E 为 \(\mathbf{C}^{2} \times \mathbf{P}^{1}(\mathbf{C})\) 的子空间，由点对 \(((u, v), x)\) 组成，这些点满足 \(v \in x\) 且 \(|u|^{2} + |v|^{2} = 1\)。在 E 上赋予 \(S_{1}\) 的作用 \(((u, v), x) \cdot z = ((zu, zv), x)\)。
\end{enumerate}

\begin{enumerate}
\def\labelenumi{\alph{enumi})}
\item
  证明第二投影 \(\mathrm{pr}_{2}\colon\mathrm{E}\to\mathbf{P}^{1}(\mathbf{C})\) 赋予 E 以 \(\mathbf{P}^{1}(\mathbf{C})\) 为底、以 \(S_{1}\) 为群的主纤维丛结构。
\item
  设 \(\varphi\) 是从 \(\mathbf{S}_2\) 到 \(\mathbf{P}^1 (\mathbf{C})\) 的同胚；证明群 \(\pi_1(\mathbf{S}_1,e)\) 由类 \(c(\varphi^{*}\mathrm{E})\) 生成。
\end{enumerate}

\backmatter

